\documentclass[10pt,a4paper]{amsart}
\usepackage[margin=19mm]{geometry}
\usepackage{amsmath,amssymb,mathtools}
\usepackage[scr=rsfs,cal=euler]{mathalfa}
\usepackage{xfrac}
\usepackage[unicode,hidelinks,bookmarks=false]{hyperref}
\newcommand{\ie}{i.e.\ }
\newcommand{\cf}{cf.\ }
\newcommand{\resp}{respectively\ }
\newcommand{\Subsection}{\subsection}
\providecommand{\nobreakdash}{\nobreak\hbox{-}}
\setlength{\emergencystretch}{2em}
\multlinegap0pt
\usepackage{tikz-cd}
\usepackage{ulem}
\newtheorem{thm}{Theorem}
\newtheorem{defn}[thm]{Definition}
\newtheorem{lem}[thm]{Lemma}
\newtheorem{prop}[thm]{Proposition}
\newtheorem{cor}[thm]{Corollary}
\newtheorem{exmp}[thm]{Example}
\newtheorem{rem}[thm]{Remark}
\newcommand{\map}[3]{{#1}:{#2}\longrightarrow{#3}}
\renewcommand{\a}{\alpha}
\renewcommand{\b}{\beta}
\newcommand{\om}{\omega}
\renewcommand{\l}{\lambda}
\newcommand{\len}[1]{\ell(#1)}
\newcommand{\N}{\mathbb{N}}
\newcommand{\pt}{\partial}
\newcommand{\last}[1]{\chi(#1)}
\newcommand{\Il}{\mathsf{L}}
\newcommand{\Ir}{\mathsf{R}}
\newcommand{\pow}[2]{{#1}^{#2}}
\newcommand{\emptyword}{\varepsilon}
\newcommand{\emptypartition}{\varepsilon}
\newcommand{\Pset}{\mathcal{P}}
\newcommand{\D}{\mathcal{D}}
\newcommand{\df}{\partial}
\newcommand{\U}{\mathcal{U}}
\newcommand{\Nil}{\mathcal{N}}
\newcommand{\cO}{\mathcal{O}}
\newcommand{\Fset}{\mathcal{F}}
\newcommand{\AP}{\mathcal{A}}
\newcommand{\BP}{\mathcal{B}}
\newcommand{\two}[1]{\mu_2(#1)}
\newcommand{\Wset}{\mathcal{W}}
\newcommand{\w}{\omega}
\newcommand{\Fbb}{\mathbb{F}}
\newcommand{\supp}{\mathcal{S}}
\newcommand{\Ba}{\mathcal{B}}
\newcommand{\Cm}{\mathcal{C}}
\newcommand{\comm}{\sim_{\mathcal{C}}}
\newcommand{\ncomm}{\not\sim_{\mathcal{C}}}
\newcommand{\Orb}{\mathcal{O}}
\newcommand{\Uc}{\mathcal{U}}
\newcommand{\Gc}{\mathcal{G}}
\newcommand{\FF}{\mathbb{F}}
\DeclareMathOperator{\GL}{GL}
\DeclareMathOperator{\im}{Im}
\DeclareMathOperator{\rank}{rank}
\newcommand{\RRset}{\mathcal{Q}}
\newcommand{\emptyf}{\emptypartition}
\newcommand{\burge}[1]{\Omega(#1)}
\newcommand{\pti}[2]{{#1}^{(#2)}}
\newcommand{\Des}[1]{\mathrm{Des}(#1)}
\newcommand{\des}[1]{\mathrm{des}(#1)}
\newcommand{\maj}[1]{\mathrm{maj}(#1)}
\newcommand{\rev}[1]{\mathrm{val}_{#1}}
\newcommand{\lev}[1]{\mathrm{val}_{#1}}
\newcommand{\ran}[1]{\mathrm{ann}_{#1}}
\newcommand{\lan}[1]{\mathrm{ann}_{#1}}
\usetikzlibrary{tikzmark}
\newcounter{pairctr}
\newcommand{\Lpair}[1]{\stepcounter{pairctr}\tikzset{tikzmark prefix=\thepairctr}\tikzmark{start}#1\tikzmark{stop}\tikz[remember picture, overlay]{
% over arrow
\draw[->]([shift={(.5ex,2.5ex)}]pic cs:start) to[bend left=20] ([shift={(-.5ex,2.5ex)}]pic cs:stop);}}
\newcommand{\Rpair}[1]{\stepcounter{pairctr}\tikzset{tikzmark prefix=\thepairctr}\tikzmark{start}#1\tikzmark{stop}\tikz[remember picture, overlay]{
% over arrow
\draw[<-]([shift={(.5ex,2.5ex)}]pic cs:start) to[bend left=20] ([shift={(-.5ex,2.5ex)}]pic cs:stop);}}
\newcommand\ara[1]{\arrow{r}{#1}}
\newcommand\arb[1]{\arrow[swap]{r}{#1}}
\newcommand\aua{\arrow{u}{\a}}
\newcommand\aub{\arrow{u}{\b}}
\newcommand\adp{\arrow[swap]{d}{\pt}}
\newcommand\auash{\arrow[u, shift right=1ex]{\a}}
\newcommand\aubsh{\arrow[u, shift right=1ex]{\b}}
\newcommand\efr{(\;)}
\newcommand{\Leila}[1]{\textcolor{purple}{[Leila] #1}}
\newcommand{\kk}{\mathsf k}
\begin{document}
\title[Jordan types commuting with a hook partition]{Jordan types commuting with a hook partition}
\author{Leila Khatami; Tomaž Košir}
\date{}
\maketitle
\noindent\emph{Source and licence.} Jordan types commuting with a hook partition. Version arXiv:2605.27692v1 (CC BY 4.0). This material is distributed under the Creative Commons Attribution 4.0 International licence, http://creativecommons.org/licenses/by/4.0/, as recorded in the arXiv OAI licence metadata for this exact version and shown on the abstract page https://arxiv.org/abs/2605.27692v1. Full rights notice: licenses/arxiv-2605.27692-CC-BY-4.0.txt.\par\smallskip
\noindent\textit{Editorial scope.} Counted unit: the original \texttt{thm} environment \texttt{all that commute with hook} at source lines 375--382 together with its complete proof at lines 384--489; the delivered document keeps source lines 156--490 so that every referenced label and every citation resolves inside it. Native editable source is the paper's own concatenated TeX; the AMS wrapper is editorial formatting only. No publisher class, upstream script or unrelated artwork is redistributed.
\section{Commuting with a hook partition}\label{hook}



In this section, we characterize all possible Jordan types in the nilpotent commutator of a hook partition $\Gamma_{n,m}=(n,1^m)$. We write $N=n+m$. Let $B$ be the nilpotent Jordan matrix with Jordan type $\Gamma_{n,m}$, and let $\U_{B}$ be the maximal nilpotent subalgebra of the nilpotent commutator $\Nil_B$ as in \cite{BaIK2010}. Then $\U_B$ intersects every nilpotent orbit that intersects $\Nil_B$ \cite[Lem. 2.2]{BaIK2010}, and consists of all matrices of the following form in $\operatorname{Mat}_{N}(\mathbb{F})$, the space of all $N$ by  $N$ matrices over the field $\mathbb{F}$.

\begin{equation}\label{generic in U}
A=\left(\begin{array}{ccccc|ccccc}
    0 & h_1 & h_2 & \cdots & h_{n-1} & u_1 & u_2 & \cdots & u_{m-1} & u_{m}\\
    0 & 0   & h_1 & \cdots & h_{n-2} & 0   & 0   & \cdots & 0        & 0\\
    \vdots &     & \ddots &        & \vdots & \vdots & \vdots &        &        & \vdots\\
    0 & 0   & 0   & \cdots & h_1     & 0   & 0   & \cdots & 0        & 0\\ 
     0 & 0   & 0   & \cdots & 0     & 0   & 0   & \cdots & 0        & 0\\
    \hline
    0 & 0   & 0   & \cdots & d_1     & 0   & 0   & 0 & \cdots       & 0\\
    0 & 0   & 0   & \cdots & d_2     & v_{12}   &0 & 0& \cdots      &0\\
    0 & 0   & 0   & \cdots & d_3     & v_{13}& v_{23}& 0 & \ddots  & 0\\
    \vdots &     &        & \cdots & \vdots & \vdots &       & \ddots & \ddots & \vdots\\
    0 & 0   & 0   & \cdots & d_m     & v_{1m} & v_{2m} & \cdots & v_{m-1,m} & 0\\
\end{array}
\right)
\end{equation}


Let $\{e_1, \ldots, e_n\}$ be the standard basis of $\mathbb{F}^n$. Then we can write 
\begin{equation}\label{blocks of generic in U}
A=\left( 
\begin{array}{c|c}
  H   & e_1 u^T \\
  \hline
  d e_n^T   & V
\end{array}
\right),
\end{equation}
where $H=\sum\limits_{i=1}^{n-1}h_iJ_n^i$, vectors $u$ and $d$ are both in $\mathbb{F}^m$, and $V$ is a strictly lower triangular matrix in $\operatorname{Mat}_{m}(\mathbb{F})$.

\begin{lem}\label{commute with hook}
    The hook partition $\Gamma_{n,m}=\left(n,1^m\right)$ commutes with any partition $Q$ of $m+n$ having one of the following forms:
    \begin{enumerate}
        \item[(a)] $Q$ is the partition obtained from $\left(\mu, [n]^k\right)$, where $\mu$ is an arbitrary partition of $m$ and $1\leq k\leq n$.
        \item[(b)] $Q$ is the partition obtained from $\left(\mu,[n-1]^k\right)$, where $\mu=(\mu_1,\ldots, \mu_s)$ is a partition of $m+1$, and $\frac{n-1}{\mu_1}\leq k \leq n-1$.
        \item[(c)] $Q$ is the partition obtained from $\left(\mu, [n-2]^k\right)$, where either $\mu=(m+2)$ and $\frac{n-1}{m+1}< k \leq n-2$, or $\mu=(\mu_1, \mu_2, \ldots, \mu_s)$ is a partition of $m+2$ such that $\mu_2\geq 2$, and $\frac{n-1}{\mu_2}< k \leq n-2$.
        \end{enumerate}
\end{lem}


\begin{proof}
Throughout this proof, we use the notations established in (\ref{generic in U}) and  (\ref{blocks of generic in U}) above.
\\

\noindent{\bf Part (a).}  Let $\mu$ be a partition of $m$, and assume that $1\leq k\leq n$. Let $Q$ be the partition obtained from $\left(\mu, [n]^k\right)$. In $\U_B$, we get Jordan type $Q$ by choosing $H=J_n^k$, $V=J_\mu^T$, and $u=d=0$ in (\ref{blocks of generic in U}) above.

\noindent{\bf Part (b).} Let $\mu=(\mu_1, \mu_2, \ldots, \mu_s)$ be a partition of $m+1$, and assume that $\frac{n-1}{\mu_1}\leq k\leq n-1$. Let $Q$ be the partition obtained from $\left(\mu, [n-1]^k\right)$.

If $\mu=\left(1^{m+1}\right)$, then $Q$ is the partition $\left(1^{m+n}\right)$, which corresponds to the zero matrix in $\U_B$.  For the rest of the proof of this part we assume that $2\leq \mu_1\leq m+1$.

Let $\mu'=(\mu_1-1,\mu_2, \ldots, \mu_s)$. Using the expression (\ref{blocks of generic in U}) for an element $A\in \U_B$, set $H=J_n^k$, $V=J_{\mu'}^T$, $u=\delta e'_{\mu_1-1}$ and $d=e_1'$, where $\{e_1', \ldots, e'_m\}$ is the standard basis of $\mathbb{F}^m$, and $\delta\in \mathbb{F}$. We reorganize the block decomposition of $A$ to get

\[A=\left( 
\begin{array}{c|c}
  J_n^k  & \delta e_1\left(e'_{\mu_1-1}\right)^T \\
  \hline
  e_1'e_n^T   & J_{\mu'}^T
\end{array}
\right)= 
\left( 
\begin{array}{c|c}
  J_{n-1}^k & f_{n-k}\left(f'_1\right)^T +\delta f_1\left(f'_{\mu_1}\right)^T\\
 \hline
  0 & J_\mu^T
\end{array}
\right),\]
%
where $\{f_1, \ldots, f_{n-1}\}$ and $\{f'_1, \ldots, f'_{m+1}\}$ are the standard bases of $\mathbb{F}^{n-1}$ and $\mathbb{F}^{m+1}$, respectively. 

Thus for a positive integer $l$, we have 

\[A^l=\left(
\begin{array}{c|c}
    J_{n-1}^{kl} & J_{n-1}^{k(l-1)}f_{n-k}(f'_1)^T+\delta f_1\left(f'_{\mu_1}\right)^T (J_\mu^{l-1})^{T}\\
    \hline
    0 & \left(J^l_{\mu }\right)^T
\end{array}
\right).
\]

\noindent {\bf Case 1.} Assume that $\frac{n}{\mu_1}\leq k$. In this case, let $\delta=0$. 

For  $1\leq l<\frac{n}{k}$, since $l<\mu_1$, the block $  \left(J^T_{\mu }\right)^l$ is non-zero and has a row equal to $(f'_1)^T$. Therefore, for such $l$, we have 
\[\rank(A^l)=\rank\left(
\begin{array}{c|c}
    J_{n-1}^{kl} &0 \\
    \hline
    0 & \left(J^l_{\mu }\right)^T
\end{array}
\right). \]

On the other hand, for $l\geq \frac{n}{k}$ we have 
\[A^l=\left(
\begin{array}{c|c}
   0 &0 \\
    \hline
    0 & \left(J^l_{\mu }\right)^T
\end{array}
\right), \]

Therefore, in this case $A$ has the desired Jordan type.

\noindent {\bf Case 2.} Now assume that $\frac{n-1}{\mu_1}\leq k <\frac{n}{\mu_1}$. This is possible only if $n-1$ is divisible by $\mu_1$ and so $n-1=k\mu_1$. Let $\delta=-1$. 

For $1\leq l \leq \mu_1-1$, we have $kl<n-1$, so that the block $J_{n-1}^{kl}$ has a column equal to $f_1$ and the block $\left(J^T_{\mu }\right)^l$ has a row equal to $(f'_1)^T$. Thus, for such $l$, we have 

\[\rank(A^l)=\rank\left(
\begin{array}{c|c}
    J_{n-1}^{kl} &0 \\
    \hline
    0 & \left(J^l_{\mu }\right)^T
\end{array}
\right). \]

For $l=\mu_1$, we have 

\[A^{\mu_1}=\left(
\begin{array}{c|c}
    0 & f_{1}(f'_1)^T- f_1\left(f'_{1}\right)^T \\
    \hline
    0 & 0
\end{array}
\right)=0=\left(
\begin{array}{c|c}
    J_{n-1}^{k\mu_1} & 0 \\
    \hline
    0 & \left(J^{\mu_1}_{\mu }\right)^{T}
\end{array}
\right), 
\]
and hence $A^l=0$ also for $l>\mu_1$.


Therefore, in this case $A$ has the desired Jordan type as well.


\noindent{\bf Part (c).} Assume that $\mu=(\mu_1,\mu_2, \ldots, \mu_s)$ is a partition of $m+2$ such that either $\mu_2=0$ or  $\mu_2\geq 2$. 

If $\mu_2=0$, let $\mu'=(m)$, and let $k$ be such that $\frac{n-1}{m+1}< k \leq n-2$. If $\mu_2\geq 2$, let $\mu'=(\mu_1-1, \mu_2-1, \mu_3, \ldots, \mu_s)$ and  $k$ be such that $\frac{n-1}{\mu_2}< k \leq n-2$.

We first note that if $\mu_1\leq 2$, then in fact $\mu_1=2$ and all parts of the partition obtained from $\left(\mu, [n-2]^k\right)$ are at most 2. We already know that such partition commutes with any partition (see \cite[Thm. 2.4]{Obl-3}). So for the rest of this proof we assume that $\mu_1\geq 3$. 

Let $A\in \U_B$, written as in (\ref{blocks of generic in U}), be such that  $H=J_n^k$, $V=\left(J_{\mu'}\right)^T$, $u=e'_{\mu_1+\mu_2-2}$, and $d=e_1'$. 
In other words, 
 \[A=\left( 
\begin{array}{c|c}
  J_n^k  & e_1\left(e'_{\mu_1+\mu_2-2}\right)^T \\
  \hline
  e_1'e_n^T   & J_{\mu'}^T
\end{array}
\right).\]

In the case $\mu=(m+2)$, for simplicity, we also write $\mu_1+\mu_2$ to mean $m+2$. 

Let $\bar{A}=P^{-1}AP$ where $P$ is the base change matrix permuting the standard basis of $\mathbb{F}^{n+m}$ by cycle  $(1\to 2\, \to \ldots \to n+\mu_1+\mu_2-2)$ and leaves other vectors fixed (if any). We can then reorganize the block decomposition of $\bar{A}$ to write 
\[\bar{A}= 
\left( 
\begin{array}{c|c}
  J_{n-2}^k & f_{n-k-1}\left(f'_1\right)^T \\
 \hline
  f'_{\mu_1+\mu_2}f_k^T & J_\mu^T
\end{array}
\right).\]

Here, $\{f_1, \ldots, f_{n-2}\}$ and $\{f_1', \ldots,f'_{m+2}\}$ are the standard bases of $\mathbb{F}^{n-2}$ and $\mathbb{F}^{m+2}$, respectively. 

Using the notation $\delta_{i,j}$ to denote the Kronecker delta, for a positive integer $l$, we have 

\[\begin{array}{ll}
\bar{A}^l&= 
\left( 
\begin{array}{c|c}
  J_{n-2}^{kl} & J_{n-2}^{k(l-1)}f_{n-k-1}\left(f'_1\right)^T \\
 \hline
  f'_{\mu_1+\mu_2}f_k^TJ_{n-2}^{k(l-1)}& \left(J^l_\mu\right)^T +\delta_{kl,n-1}f'_{\mu_1+\mu_2}(f'_1)^T
\end{array}
\right)
\\ \\
&=
\left( 
\begin{array}{c|c}
  J_{n-2}^{kl} &f_{n-kl-1}\left(f'_1\right)^T \\
 \hline
  f'_{\mu_1+\mu_2}f_{kl}^T& \left(J^l_\mu\right)^T +\delta_{kl,n-1}f'_{\mu_1+\mu_2}(f'_1)^T
\end{array}
\right).
\end{array}\]

We will show that $A$ has the desired Jordan type by showing that for every positive integer $l$, 
\[\rank\bar{A}^l=\rank \left( 
\begin{array}{c|c}
  J_{n-2}^{kl} & 0 \\
 \hline
  0& \left(J^l_\mu\right)^T
\end{array}
\right).\]

\noindent {\bf Case I.} Assume that $\mu_2\geq 2$. If $l$ is such that $1\leq l <\mu_2$, then the $(\mu_1+\mu_2-l)$-th column of $\left(J^l_\mu\right)^T$ is equal to $f'_{\mu_1+\mu_2}$. Therefore, for such $l$, we can use an elementary column operation to remove the non-zero entry in the block $\bar{A}^l_{21}$, which will be in row $\mu_1+\mu_2$ and column $kl$ of the block, as well as the extra term in the block $\bar{A}^l_{22}$, which will be in row $\mu_1+\mu_2$ and column $1$ of the block. Moreover, the $(l+1)$-st row of $\left(J^l_\mu\right)^T$ is equal to $(f'_1)^T$ and therefore a nonzero entry in the block $\bar{A}^l_{12}$, which will be in row $n-kl-1$ and column $1$ of that block, can be eliminated using an elementary row operation. 

On the other hand, for $l\geq \mu_2$, since $kl> n-1$, both blocks $\bar{A}^l_{12}$ and $\bar{A}^l_{21}$, as well as the extra term in the block $\bar{A}^l_{22}$, are 0.
\\

\noindent {\bf Case II.} Now, assume that $\mu=(m+2)$ and that $\frac{n-1}{m+1}< k \leq n-2$. We note that in this case, as long as $l\leq m$, the $(l+1)$-st row of $\left(J^l_\mu\right)^T$ is $(f'_1)^T$, and its $(m+2-l)$-th column is $f'_{m+2}$. Since $l+1\leq m+1$ and $m+2-l\geq 2$, a nonzero entry in $\bar{A}^l_{12}$ or $\bar{A}^l_{21}$, or an additional term in $\bar{A}^l_{22}$ can be eliminated through appropriate elementary column and row operations.

On the other hand, for $l\geq m+1$, we have $lk> n-1$ and therefore both $\bar{A}^l_{12}$ and $\bar{A}^l_{21}$, as well as the extra term in $\bar{A}^l_{22}$ are 0.

Therefore, in both cases, for every positive integer $l$, we have the desired rank conditions, implying that $A$ has Jordan type $([n-2]^k, \mu)$.
\end{proof}

\begin{rem}
In part (c) of Lemma~(\ref{commute with hook}), we made the lower bounds for $k$ sharp, but partitions that come from the inclusive lower bounds for $k$ are also among the commuting partitions covered in earlier parts of Lemma \ref{commute with hook}. If $\mu=(m+2)$ and $k=\frac{n-1}{m+1}$, then $([n-2]^k, m+2)=(m,(m+1)^{k-1},m+2)=([n]^k, m)$, which is already covered in Part (a). On the other hand, if $\mu=(\mu_1, \mu_2, \ldots, \mu_s)$ with $\mu_2\geq 2$, for $k=\frac{n-1}{\mu_2}$, we have $([n-2]^k, \mu)=(\mu_2^{k-1}, \mu_2-1, \mu)=([n-1]^k, \mu_1,\mu_2-1, \ldots, \mu_s)$, which is already covered in case (b) because $\frac{n-1}{\mu_1}\leq \frac{n-1}{\mu_2}\leq k$.
\end{rem}


\begin{thm}\label{all that commute with hook}
    Partition $\Gamma_{n,m} =\left(n,1^m\right)$ commutes with a partition $Q$ of $m+n$ if and only if $Q$ is one of the following partitions:
    \begin{enumerate}
        \item[(a)] $Q$ is the partition obtained from $\left(\mu, [n]^k\right)$, where $\mu$ is an arbitrary partition of $m$ and $1\leq k\leq n$.
        \item[(b)] $Q$ is the partition obtained from $\left(\mu,[n-1]^k\right)$, where $\mu=(\mu_1,\ldots, \mu_s)$ is a partition of $m+1$, and $\frac{n-1}{\mu_1}\leq k\leq n-1$.
        \item[(c)] $Q$ is the partition obtained from $\left(\mu, [n-2]^k\right)$, where either $\mu=(m+2)$ and $\frac{n-1}{m+1}< k \leq n-2$, or $\mu=(\mu_1, \mu_2, \ldots, \mu_s)$ is a partition of $m+2$ such that $\mu_2\geq 2$, and $\frac{n-1}{\mu_2} < k \leq n-2$.
    \end{enumerate}
\end{thm}

\begin{proof}
A generic matrix $A$ in $\U_B$ as in \eqref{generic in U} is written in the block form as in \eqref{blocks of generic in U}. We know that $H$ is similar to $J_n^k$ for some power $k\in\{1,2,\ldots,n\} $ and $V$ is similar to a nilpotent Jordan matrix with Jordan type $\nu=(\nu_1,\nu_2,\ldots,\nu_t)$. Then we have $H=S^{-1} J_n^k S $ and $V=T^{-1}J_{\nu}T $ for two invertible matrices $S$ and $T$. Since $H$ is strictly upper-triangular matrix we may choose $S$ such that $Se_1=\alpha e_1$ and $S^Te_n=\beta e_n$ for two nonzero scalars $\alpha$ and $\beta$. Then, we have
$$
A=\left( 
\begin{array}{c|c}
  S^{-1} J_n^kS   & e_1 u^T \\
  \hline
  d e_n^T   & T^{-1}J_{\nu}T
\end{array}
\right)=\left( 
\begin{array}{c|c}
  S^{-1}   & 0 \\
\hline
0   & T^{-1}
\end{array}
\right)\left( 
\begin{array}{c|c}
  J_n^k   & e_1 a^T \\
  \hline
  b e_n^T   & J_{\nu}
\end{array}
\right)\left( 
\begin{array}{c|c}
  S   & 0 \\
  \hline
  0   & T
\end{array}
\right),
$$
where we write $u=\alpha^{-1}T^Ta $ and $d=\beta T^{-1}b $.
Thus $A$ is similar to 
\[B=\left( 
\begin{array}{c|c}
  J_n^k   & e_1 a^T \\
  \hline
  b e_n^T   & J_{\nu}
\end{array}
\right)\]
and we will consider $B$ hereafter.

In order to find the Jordan type of $B$ we will compute its powers. Inductively, we prove that
$$
B^l=\left( 
\begin{array}{c|c}
  J_n^{kl} + \left(a^TJ_{\nu}^{l-2}b\right) J_n^{n-1}   & e_1 a^T J_{\nu}^{l-1} \\
  \hline
   J_{\nu}^{l-1} b e^T_n  & J_{\nu}^l
\end{array}
\right)
$$
for any $l\ge 2$. As long as $kl\leq n-2$, the off diagonal blocks do not affect the rank of $B^l$ and we have $\rank(B^l)=\rank(J_n^{kl})+\rank(J_{\nu}^l )$.

Consider next the case $kl=n-1$. We omit the first $n-1$ consecutive columns and $n-1$ consecutive rows, starting with the second row, in $B^l$ that are all zero to obtain
\begin{equation}
\label{kl=n-1}
\left( 
\begin{array}{c|c}
  1 + \left(a^TJ_{\nu}^{l-2}b\right)    & a^T J_{\nu}^{l-1} \\
  \hline
   J_{\nu}^{l-1} b   & J_{\nu}^l
\end{array}
\right).
\end{equation}
We consider several cases depending on the largest part $\nu_1$ of partition $\nu$. 

{\bf Case I.} $\nu_1\leq l-2$.

In this case we have $J_{\nu}^{l-2}=0$, and thus $\rank B^{l}=1$. We also have $\rank B^{s}=0$ for all $s>l$. The Jordan type of $B$ is equal to 
$\left(\nu, [n]^k\right)$ where $\nu$ is a partition of $m$ and $1\leq k\leq n$. So it belongs to the case (a) of Lemma \ref{commute with hook}.

{\bf Case II.} $\nu_1= l-1.$
Now, we have $J_{\nu}^{l-1} =J_{\nu} ^l=0$. So, $\rank B^l$ is either $0$ or $1$, depending on $1+a^TJ_{\nu}^{l-2}b$ being equal to $0$ or not, respectively. Suppose first that $\rank B^l=1$. Then, it follows that $\rank B^{s}=0$ for all $s>l$. Hence $\rank(B^l)=\rank(J_n^l)+\rank(J_{\nu}^l )$ for all $l$ and the Jordan type of $B$ is equal to 
$\left(\nu, [n]^k\right)$ where $\nu$ is a partition of $m$ and $1\leq k\leq n$. The Jordan type of $A$ in this subcase belongs to the case (a) of Lemma \ref{commute with hook}.

Suppose next that $\rank B^l=0$. Since $\nu_1=l-1$ and $kl=n-1$ it follows that the Jordan type of $B$ is equal to $\left(\nu', [n-1]^k\right)$, where $\nu'=(\nu_1+1,\nu_2,\nu_3,\ldots, \nu_t)=(l,\nu_2,\ldots,\nu_t)$ and $[n-1]^k=(l^k)$. Since $\nu_1+1=l=\frac{n-1}{k}$, it follows that the Jordan type in this case belongs to type (b) of Lemma \ref{commute with hook}.

{\bf Case III.} $\nu_1\ge l.$

Now the rank of $B^l$ can be equal to one of $3$ possible values: $\rank J_{\nu}^l$, $\rank J_{\nu}^l+1$ or $\rank J_{\nu}^l+2$, depending on $a$ and $b$. We consider the three cases separately.

\begin{itemize}
    \item[(i)] $\rank B^l-\rank J_{\nu}^l=0$: Then $1+a^TJ_{\nu}^{l-2}b=0 $, $J_{\nu}^{l-1}b\in\im J_{\nu}^l $ and $\left(J_{\nu}^{l-1}\right)^Ta\in\im \left(J_{\nu}^l\right)^T $. We have $\rank B^j=n-jk+\rank J_{\nu}^j $ for $j=0,1,\ldots,l-1$, and $\rank B^l=\rank J_{\nu}^l $. Since $n=kl+1$ it follows that the Jordan type of $B$ is equal to $([n-1]^k,\nu' ) $, where $\nu'$ is a partition of $m+1$ and $\nu_1'\ge\nu_1\ge l  = \frac{n-1}{k} $. 

So the Jordan type of $B$ belongs to type (b) of Lemma \ref{commute with hook}.

\item[(ii)] $\rank B^l-\rank J_{\nu}^l=1$: Then $\rank B^j=\rank J_n^{kj}+\rank J_{\nu}^j $ for all $j\leq l$. Since $J_n^{kl+1}=J_n^n=0$, the equality $\rank B^j=\rank J_n^{kj}+\rank J_{\nu}^j$ holds also for all $j> l$.  So, the Jordan type of $B$ is equal to $\left(\nu', [n]^k\right)$ where $\nu'$ is a partition of $m$ and $1\leq k\leq n$. It belongs to the case (a) of Lemma \ref{commute with hook}.

\item[(iii)] $\rank B^l-\rank J_{\nu}^l=2$: Then $\rank B^j=\rank J_n^{kj}+\rank J_{\nu}^j $ for all $j\leq l-1$ and $\rank B^l=\rank J_{\nu}^l+2=\rank J_n^{n-1}+\rank J_{\nu}^l+1$. Therefore, the Jordan type of $B$ is equal to $([n-1]^k,\nu')$, where $\nu'=(\nu'_1,\nu'_2,\ldots,\nu'_t)$ is a partition of $m+1$ and $\nu_1'\ge\nu_1+1\ge l+1\ge \frac{n-1}{k} $. It belongs to the case (b) of Lemma \ref{commute with hook}.

\end{itemize}

It remains to consider the case when there is no $l$ such that $kl=n-1$, i.e., we assume that the remainder of $n$ when divided by $k$ is not equal to $1$. We write $\ell=\left\lfloor{\frac{n}{k}}\right\rfloor$. Then 
$\rank B^j=\rank J_n^{kj}+\rank J_{\nu}^j $ for all $j\leq \ell$. 

\begin{itemize}
    \item[(i)] If $\nu_1\leq\ell$ then $B^j=0$ for $j\ge\ell +1$ and the Jordan type of $B$ is equal to $\left([n]^k,\nu\right)$ where $\nu$ is a partition of $m$ and $1\leq k\leq n$. It belongs to the case (a) of Lemma \ref{commute with hook}.

\item[(ii)] Assume next that $\nu_1\ge\ell+1$, and that $a^TJ_{\nu}^{\ell-1}b=0 $, $J_{\nu}^{\ell}b\in\im J_{\nu}^{\ell +1} $ and $\left(J_{\nu}^{\ell}\right)^Ta\in\im J_{\nu}^{\ell +1} $. Then, we have $\rank B^j=\rank J_n^{kj}+\rank J_{\nu}^j$  for all $j$ and the Jordan type of $B$ is equal to $\left([n]^k,\nu\right)$ where $\nu$ is a partition of $m$ and $1\leq k\leq n$. It belongs to the case (a) of Lemma \ref{commute with hook}.

\item[(iii)] Assume next that $\nu_1\ge\ell+1$, and that $a^TJ_{\nu}^{\ell-1}b\neq 0 $, but $J_{\nu}^{\ell}b\in\im J_{\nu}^{\ell+1}$ and $\left(J_{\nu}^{\ell}\right)^Ta\in\im J_{\nu}^{\ell+1}  $. Then $\rank B^j=\rank J_n^{kj}+\rank J_{\nu}^j $ for all $j\leq \ell$ and $\rank B^{\ell +1} =\rank J_{\nu}^{\ell +1} +1$. Therefore, the Jordan type of $B$ is equal to $([n-1]^k,\nu')$, where $\nu'=(\nu'_1,\nu'_2,\ldots,\nu'_t)$ is a partition of $m+1$ and $\nu_1'\ge\nu_1\ge \ell+1\ge \frac{n-1}{k} $. It belongs to the case (b) of Lemma \ref{commute with hook}.

\item[(iv)] If $\nu_1\ge\ell + 1$ but exactly one of the relations $J_{\nu}^{\ell}b\in\im J_{\nu}^{\ell +1}$ and $\left(J_{\nu}^{\ell}\right)^Ta\in\im J_{\nu}^{\ell +1} $ holds and the other fails, then $\rank B^j=\rank J_n^{kj}+\rank J_{\nu}^j $ for all $j\leq \ell$ and $\rank B^{\ell +1} =\rank J_{\nu}^{\ell +1} +1$. Therefore, the Jordan type of $B$ is equal to $([n-1]^k,\nu')$, where $\nu'=(\nu'_1,\nu'_2,\ldots,\nu'_t)$ is a partition of $m+1$ and $\nu_1'\ge\nu_1\ge \ell+1\ge \frac{n-1}{k} $. It belongs to the case (b) of Lemma \ref{commute with hook}.

\item[(v)] Finally, assume that $\nu_1\ge\ell +1$ and both relations $J_{\nu}^{\ell}b\notin\im J_{\nu}^{\ell +1}$ and $\left(J_{\nu}^{\ell}\right)^Ta\notin\im J_{\nu}^{\ell +1} $ hold. Then $\rank B^j=\rank J_n^{kj}+\rank J_{\nu}^j $ for all $j\leq \ell$ and $\rank B^{\ell +1} =\rank J_{\nu}^{\ell +1} +2$. Therefore, the Jordan type of $B$ is equal to $([n-2]^k,\nu')$, where $\nu'=(\nu_1,\nu_2,\ldots,\nu_s)$ is a partition of $m+2$. If $\nu=(m)$ then $s=1$ and $\nu'=(m+2)$. If $\nu_2\ge 1$ then $\nu'$ is a partition of $m+2$ such that $\nu_2'\ge \ell+1\ge 2$ and $\frac{n-1}{\nu'_2}\leq k \leq n-2$. This case belongs to the case (b) or (c) of Lemma \ref{commute with hook}.
\end{itemize}
\end{proof}


\begin{thebibliography}{99}
\bibitem[BaIK2010]{BaIK2010} R. Basili, A. Iarrobino, L. Khatami. \emph{Commuting nilpotent matrices and Artinian algebras.} J. Commut. Alg. \textbf{2} (2010), no. 3, 295-325.
\bibitem[Obl-3]{Obl-3} P. Oblak. \emph{On the nilpotent commutator of a nilpotent matrix.} Lin. Multilin. Algebra \textbf{60} (2012), no. 5, 599-612.
\end{thebibliography}
\end{document}
